% cap1.tex
\chapter{Introduction}
\label{ch:intro} % la etiqueta para referencias

\color{blue}{En esta sección de debe incluir el fundamento general de la motivación de la investigación desarrollada en la tesis. También se debe incluir la Hipótesis de trabajo, Objetivo General y específicos. (Máximo 5 páginas). Usar mismo tamaño, tipo de letra y espacio del este formato).}
\vspace{0.5cm}



\section{Hypothesis}


\color{blue}{The hypothesis addressed by this thesis is divided into the following research questions:}

\begin{itemize}
   \item RQ1: 
   \item RQ2: 
 \end{itemize}






\section{Objetives}

\subsection{General Objetive}

\vspace{0.5cm}


\vspace{0.5cm}
\subsection{Specific Objectives }

\begin{enumerate}[]
\item 
\item 
\item 
\end{enumerate}





\section{Thesis Structure}

% {\color{blue} (2 a 3 líneas que informe al lector el contenido de cada capítulo) }

% \textcolor{red}{ALONSO ARMAR LA ESTRUCTURA DE ESTA FORMA Y AUTOCOMPLETAR !!!!
% This document is organized into five chapters, which are described below:}

\color{blue}{This thesis is organized into five distinct chapters, each contributing uniquely to the overarching research endeavor:}
\begin{enumerate}
\item  Chapter \ref{ch:intro}: 
\item  Chapter \ref{ch:litRev}: 
\item Chapter \ref{ch:metodo}: 
\item Chapter \ref{ch:resultados}: 
\item Chapter \ref{ch:conclu}: 
\end{enumerate}

% Following this structure, the thesis aims to contribute to the development of an \textcolor{red}{ALONSO AUTOCOMPLETAR !!!!}

